Let be $R$ the radius of the Celestial Sphere and $(x, y, z)$ a coordinate
system with origin at the observer, $z$ is the down-up axis, $y$ is the
east-west axis and $x$ is the south-north axis.

In the superior culmination, the altitude angle is
$h = 7.41 + 49.35 = 56.76^\circ$ and we have
\begin{align*}
  \vec{v} &= v_0\,\hat{y} = 2\pi R\,\hat{y}\ /day \\
  \vec{r}_0 &= R(-\cos h\,\hat{x} + \sin h\,\hat{z}) \\
  \mathbf{r}(t) &= \mathbf{r}_0 + \mathbf{v}t
\end{align*}

The altitude in function of our coordinate system is
$h = \tan^{-1}\left(\frac{z}{\sqrt{x^2 + y^2}}\right)$. When $t \to \infty$,
the coordinates $x$ and $z$ of the star remain constant, while
$y \to \infty$, and we have
\[
  h = \tan^{-1}\left(\frac{z}{\sqrt{x^2 + y^2}}\right)
    \xrightarrow[t \to \infty]{} \tan^{-1} 0 = 0^\circ
\]
\hfill(3 pt $h \to 0$)

Similarly, because $y \to \infty$ but $x$ and $z$ remain constant, the final
azimuth is the West direction ($A = 90^\circ$ or $A = 270^\circ$, depending on
the system used).
\hfill(2 pt $A \to$ W)

For the magnitude,
$m_0 - m_6 = -2.5\log\frac{F_0}{F_6}$ and $F = \frac{L}{4\pi R^2}$ so
$R_6 = 12.88\,R_0$
\hfill(2 pt ratio of distances)
\[
  |\mathbf{r}(t)| = \sqrt{x^2 + y^2 + z^2} = R\sqrt{1 + 4\pi^2 t^2} = 12.88\,R
\]
\[
  t = 2.044\ \text{days}
\]
\hfill(3 pt time)
